% !TeX spellcheck = ru_RU
% !TEX root = report.tex

\section{Обзор}
\label{sec:review}

В разделе изложены определения и алгоритмы статьи~\cite{KosovskayaIsom}, на которых основана работа, а также известные подходы к проверке изоморфизма графов.

\subsection{Гиперграфы отношений}

Пусть на множестве $\omega = \{\omega_1, \dots, \omega_t\}$ заданы свойства и отношения $p_1, \dots, p_n$.
Гиперграф отношений~\cite{KosovskayaHypergraphs2025} имеет \emph{предикатные} вершины, помеченные именами отношений, и \emph{объектные}, помеченные элементами~$\omega$.
Ребро типа~(1) $p^k \to (\omega_{i_1}, \dots, \omega_{i_k})$ соединяет вершину некоммутативного $k$-местного отношения с упорядоченным набором различных объектных вершин, ребро типа~(2) $p^{c,k} \to \{\omega_{i_1}, \dots, \omega_{i_k}\}$~--- вершину коммутативного отношения с их множеством.
Так, пятеро братьев описываются одним ребром $\mathit{brothers}^{c,5} \to \{\omega_1, \dots, \omega_5\}$.

\begin{definition}[\cite{KosovskayaIsom}]
    Гиперграфы отношений $G_1$ и $G_2$ изоморфны, если существуют гиперграф $H$, у которого вместо имён объектных вершин стоят переменные, и подстановки (\emph{унификаторы}) $\lambda_1$, $\lambda_2$ имён объектных вершин вместо переменных такие, что $H\lambda_1 = G_1$ и $H\lambda_2 = G_2$.
\end{definition}

На практике $H$ получают из $G_2$, заменяя различные константы $b_1, \dots, b_n$ различными переменными $x_1, \dots, x_n$, и ищут биекцию $\lambda$ переменных $H$ на объектные вершины $G_1$, для которой $H\lambda = G_1$.
Задача проверки изоморфизма таких гиперграфов полиномиально эквивалентна задаче изоморфизма графов~\cite{KosovskayaIsom, KosovskayaKosovskii2019}.

\subsection{Характеристики и структурное отображение}

Пусть гиперграф содержит одну предикатную вершину.
\emph{Характеристика} вершины $a$ для рёбер типа~(1)~--- вектор $\chi(a) = (n_1, \dots, n_k)$, где $n_j$~--- число вхождений $a$ на $j$-е место; для рёбер типа~(2)~--- число вхождений $a$ в рёбра.
\emph{Характеристика гиперграфа} $\chi(G)$~--- упорядоченный набор характеристик его вершин; для проверки равенства годится любой канонический порядок, например лексикографический.

\begin{theorem}[\cite{KosovskayaIsom}]
    Для изоморфизма гиперграфов с одной предикатной вершиной необходимо, чтобы их характеристики совпадали.
\end{theorem}

Вершины с одинаковыми характеристиками образуют классы.
\emph{Структурное отображение} $\{[x_{i_1}, \dots, x_{i_{t_1}}] : [a_{j_1}, \dots, a_{j_{t_1}}], \dots\}$ сопоставляет каждому классу переменных $H$ класс констант $G$ с той же характеристикой.
Переменные класса могут переходить только в константы того же класса, поэтому вместо $n!$ подстановок достаточно перебрать $\prod_{k=1}^{K} t_k!$, где $t_1, \dots, t_K$~--- размеры классов.
Например, в первом примере статьи (отношение $p^3$, шесть вершин) отображение имеет вид $\{[x_5]:[a_4],\ [x_3]:[a_3],\ [x_2, x_6]:[a_1, a_5],\ [x_1, x_4]:[a_2, a_6]\}$, и вместо $6! = 720$ подстановок остаётся~4.

\subsection{Алгоритмы \ISOMone{} и \ISOMn{}}

Алгоритм \ISOMone{} проверяет изоморфизм гиперграфов $G_1$ и $H$ с одной предикатной вершиной.
Он получен перестройкой алгоритма для предикатных формул с одним предикатным символом~\cite{ZhouKosovskaya2024} на гиперграфы с рёбрами типов~(1) и~(2).
Алгоритм сравнивает характеристики (шаги~1--2), строит структурное отображение (шаг~3) и перебирает совместимые с ним подстановки $\lambda$, проверяя $H\lambda = G_1$ (шаги~4--7).
Подстановку отбрасывают досрочно, если часть $H\lambda$, содержащая уже известную константу, не согласуется с соответствующей частью $G_1$ (шаги~8--12).
При известном отображении перебирается не более $\prod_{k=1}^{K} t_k!$ возможных подстановок групп констант вместо групп переменных, задаваемых отображением.
Алгоритм без шагов~1--3, получающий готовое отображение, называется \ISOMm{}.

Пусть теперь предикатных вершин несколько, и $G^i$, $H^i$~--- подграфы, содержащие только $i$-ю из них.
Сначала сравниваются характеристики всех пар $G^i$, $H^i$.
Затем алгоритм \ISOMn{} строит дерево частичных унификаторов: тройки $\langle G, H, \lambda \rangle$ хранятся в очереди, упорядоченной по числу возможных подстановок $\prod t_k!$, и обрабатываются так (рис.~\ref{fig:isomn}):
\begin{enumerate}
    \item \proc{MAP}: построить отображения $\mathit{mapping}^i$ для подграфов, в которых остались переменные;
    \item \proc{BESTMAP}: выбрать подграф $i_0$ с наименьшим числом возможных подстановок $\prod t_k!$;
    \item \ISOMm{}: найти локальные унификаторы $\hat\lambda_1^{i_0}, \dots, \hat\lambda_r^{i_0}$ подграфов $G^{i_0}$ и $H^{i_0}\lambda$;
    \item \proc{CONSISTENCY}: оставить те из них, которые вместе с $\lambda$ образуют биекцию (нет пар $x/a,\ x/b$ и $x/a,\ y/a$ при $x \ne y$, $a \ne b$);
    \item \proc{REDUCTION}: для каждого оставшегося $\lambda_t$ сравнить с $G^j$ все подграфы $H^j\lambda\lambda_t$ без переменных; при несовпадении $\lambda_t$ отбрасывается, а если переменных не осталось нигде, $\lambda\lambda_t$~--- унификатор, иначе тройка $\langle G, H, \lambda\lambda_t \rangle$ ставится в очередь.
\end{enumerate}

\begin{figure}[ht]
    \centering
    \begin{tikzpicture}[
            font=\small, node distance=4mm and 9mm,
            block/.style={draw, rectangle, minimum width=34mm, minimum height=6.5mm, align=center},
            term/.style={draw, rounded rectangle, minimum height=6.5mm, align=center},
            test/.style={draw, diamond, aspect=2.4, inner sep=1pt, align=center},
            arr/.style={-{Stealth[length=2mm]}, line width=0.5pt},
        ]
        \node[term] (start) {$\langle G, H, \varnothing \rangle$};
        \node[block, below=of start] (map) {\proc{MAP}};
        \node[block, below=of map] (best) {\proc{BESTMAP}: $i_0$, $\min \prod t_k!$};
        \node[block, below=of best] (isom) {\ISOMm{}$(G^{i_0}, H^{i_0}\lambda)$};
        \node[test, below=of isom] (r0) {$r = 0$?};
        \node[block, below=of r0] (cons) {\proc{CONSISTENCY}};
        \node[test, below=of cons] (t0) {$T = 0$?};
        \node[block, below=of t0] (red) {\proc{REDUCTION}};
        \node[test, below=of red] (empty) {очередь\\ пуста?};
        \node[term, below=of empty] (stop) {конец};
        \node[term, right=12mm of empty] (pop) {$\langle G, H, \lambda \rangle$\\ из очереди};
        \node[term, left=10mm of red] (out) {полные~--- в ответ,\\ частичные~--- в очередь};

        \draw[arr] (start) -- (map);
        \draw[arr] (map) -- (best);
        \draw[arr] (best) -- (isom);
        \draw[arr] (isom) -- (r0);
        \draw[arr] (r0) -- node[right] {нет} (cons);
        \draw[arr] (cons) -- (t0);
        \draw[arr] (t0) -- node[right] {нет} (red);
        \draw[arr] (red) -- (empty);
        \draw[arr] (empty) -- node[right] {да} (stop);
        \draw[arr] (empty) -- node[above] {нет} (pop);
        \draw[arr] (pop) |- (map);
        \draw[arr, densely dotted] (red) -- (out);
        \coordinate (left) at ($(out.west) + (-4mm, 0)$);
        \draw[arr] (r0.west) node[above left] {да} -| ($(left) + (-3mm, 0)$) |- (empty.west);
        \draw[line width=0.5pt] (t0.west) node[above left] {да} -| ($(left) + (-3mm, 0)$);
    \end{tikzpicture}
    \caption{Блок-схема алгоритма \ISOMn{} (по рис.~3 из~\cite{KosovskayaIsom})}
    \label{fig:isomn}
\end{figure}

\subsection{Пример работы \ISOMn{}}
\label{subsec:example}

Рассмотрим гиперграфы
\begin{align*}
    G_1 & = \{p_1 \to a_1,\ p_1 \to a_2,\ p_1 \to a_3,\ p_2^2 \to (a_1, a_2),\ p_2^2 \to (a_1, a_4), \\
        & \phantom{{}=\{} p_3^{c,3} \to \{a_2, a_3, a_4\}\},                                          \\
    G_2 & = \{p_1 \to b_2,\ p_1 \to b_3,\ p_1 \to b_4,\ p_2^2 \to (b_4, b_3),\ p_2^2 \to (b_4, b_1), \\
        & \phantom{{}=\{} p_3^{c,3} \to \{b_1, b_2, b_3\}\}
\end{align*}
и шаблон $H$, полученный из $G_2$ заменой $b_i$ на $x_i$.
\proc{BESTMAP} выбирает подграф $p_2^2$ с числом возможных подстановок $1! \cdot 2! = 2$; из двух его локальных унификаторов один отсекается на шаге \proc{MAP}, второй после обработки $p_1$ даёт унификатор (рис.~\ref{fig:tree}).

\begin{figure}[ht]
    \centering
    \begin{tikzpicture}[
            font=\small, level distance=20mm, sibling distance=78mm,
            every node/.style={align=center},
            edge from parent/.style={draw, -{Stealth[length=2mm]}, line width=0.5pt},
        ]
        \node {$\lambda_0 = \varnothing$\\ \proc{BESTMAP}: $p_2^2$, $\{[x_1, x_3]:[a_2, a_4],\ [x_4]:[a_1]\}$}
        child {node {$x_1 \mapsto a_2,\ x_3 \mapsto a_4,\ x_4 \mapsto a_1$\\ \proc{MAP}: $a_4$ не входит в $G^{p_1}$\\ \textbf{ветвь отсечена}}}
        child {node {$x_1 \mapsto a_4,\ x_3 \mapsto a_2,\ x_4 \mapsto a_1$\\ \proc{BESTMAP}: $p_1$, $\{[x_2]:[a_3]\}$}
                child {node {$x_2 \mapsto a_3$\\ \textbf{унификатор}}}};
    \end{tikzpicture}
    \caption{Дерево поиска \ISOMn{} для примера из раздела~\ref{subsec:example}}
    \label{fig:tree}
\end{figure}

\subsection{Существующие подходы и выбор средств реализации}

Задача изоморфизма графов не доказана ни полиномиальной, ни NP-полной; лучший известный алгоритм квазиполиномиален~\cite{Babai2016}.
На практике используют программы канонической разметки, такие как nauty и Traces~\cite{McKayPiperno2014}, основанные на итеративном уточнении разбиения вершин, и перебор с отсечениями, например VF2~\cite{Cordella2004}.
Изоморфизм гиперграфов сводится к изоморфизму двудольного графа инцидентности~\cite{Luks1999}; для рёбер типа~(1) понадобились бы дополнительные вершины для номеров мест.
Подход статьи~\cite{KosovskayaIsom} работает с гиперграфами отношений непосредственно: характеристики делят вершины на классы, перебор идёт по перестановкам внутри классов, а порядок обработки подграфов выбирается по числу возможных подстановок.

Python позволяет записать алгоритм почти дословно по статье: подстановка~--- словарь, мультимножество рёбер~--- \code{Counter}, перебор перестановок~--- \code{itertools.permutations}; но интерпретатор медленный.
\Cpp{}17 позволяет хранить данные в компактных массивах, а нужные средства (\code{std::next\_permutation}, кучи) есть в стандартной библиотеке.

\subsection{Выводы}

Алгоритм \ISOMn{} сочетает необходимое условие изоморфизма, сокращение перебора за счёт классов вершин и поиск по дереву частичных унификаторов с выбором подграфа с наименьшим числом возможных подстановок.
Число возможных подстановок определяется размерами классов, поэтому время работы должно сильно зависеть от структуры входа; это нужно учесть при выборе тестовых задач.
Для проверки реализаций нужен независимый эталон; в работе им служит полный перебор всех $n!$ биекций.
